\documentclass[conference]{IEEEtran}
\usepackage{array,tabularx}
\usepackage{graphicx,xcolor,float}
\usepackage[caption=false,font=footnotesize]{subfig}
\usepackage[hidelinks]{hyperref}

% Fill in the placeholders. Add table rows as needed.
\renewcommand{\thesection}{\arabic{section}}
\newcolumntype{Y}{>{\raggedright\arraybackslash}X}
\renewcommand{\arraystretch}{1.25}
\newcommand{\fillin}[1]{\textcolor{gray}{\textit{[#1]}}}
\newcommand{\entry}[2]{\par\medskip\noindent\textbf{#1}\par\fillin{#2}\par}
\newcommand{\blank}{\fillin{---}}
\newcommand{\figureplaceholder}[1]{\fbox{\parbox[c][1in][c]{0.92\linewidth}{\centering\fillin{#1}}}}

\title{Lab 4: Design, Prototyping, and Anthropomorphization\\
\large INFO 5356: Introduction to Human-Robot Interaction}
\author{
\IEEEauthorblockN{\fillin{Elliot Zhang, Yifan (Fanning) Luo, Jaclyn Pham, Jiaxin (Aurora) Shen}}
\IEEEauthorblockA{Cornell Tech\\
\{\fillin{qz432, yl4389, cqp4, js3996}\}@cornell.edu}
}

\begin{document}
\maketitle


\section{Introduction}

Disco Moo is an educational human-robot interaction game designed to introduce K--12 students, particularly middle school learners, to foundational programming concepts through tangible, block-based coding. Inspired by Scratch, the game transforms conventional screen-based programming into a physical activity in which students assemble 3D-printed coding blocks to control a cow-themed Reachy Mini robot.

Each block represents a programming instruction, such as turning the robot's head, performing a movement, or repeating an action. The blocks incorporate ArUco markers that allow the system to visually recognize and interpret the assembled commands. During gameplay, Reachy presents a dance challenge, and students construct a sequence of blocks to produce the requested movements. Through a camera-based interface, the program can be scanned and executed by Reachy, providing a physical demonstration of the relationship between code and robotic behavior.

The game primarily focuses on three computational thinking concepts: sequencing, repetition, and debugging. By arranging, testing, and revising physical instructions, students can explore how programming logic translates into observable actions without needing prior experience writing Python code. The design also creates opportunities for collaborative problem-solving and progressive learning, informed by feedback from a K--12 teacher panel.

For Lab 4, our team developed a preliminary functional prototype, including 3D-printed interlocking coding blocks, paper-based ArUco marker recognition, a camera-based execution interface, and a customized cow-inspired robot embodiment. These components establish the foundation for further integration and evaluation in the final project.



\section{Teacher Interviews}

\subsection{Interview Questions and Responses}

Our team participated in the K--12 Teacher Panel on October 28, 2026, to gather educators' perspectives on designing an interactive robotic game for middle school students. We collected notes from five teachers.

We prepared four questions to understand students' programming experiences, common learning challenges, and opportunities for integrating physical robotics into computer science education:

\begin{enumerate}
    \item What CS courses or topics do you teach to middle school students, and at which grade levels?
    \item Which programming languages, platforms, or tools do your students use?
    \item Which programming concepts do you teach, and which do students find most challenging?
    \item How could physical coding blocks and a moving robot support the concepts you teach? Which concepts seem like the best fit, and why?
\end{enumerate}

The teachers discussed Scratch, block-based programming, and Python as common tools for introducing programming concepts. They emphasized that students should develop foundational computational thinking skills rather than rely solely on AI-assisted coding tools.

Teachers also highlighted the importance of breaking complex programming tasks into manageable steps. Visual and verbal instructions were recommended to support students with different learning needs, including English language learners and students with disabilities.

Regarding feedback, teachers recommended providing specific, encouraging hints instead of immediately revealing correct answers. They also emphasized collaborative activities and rotating student roles to ensure participation when classroom access to robots is limited.

Finally, teachers stressed that educational robotics activities should align with curriculum objectives and remain feasible within classroom time and resource constraints.

\subsection{Thematic Analysis}

We conducted an inductive thematic analysis of the October 28 teacher panel transcript. Each team member identified excerpts that bore on the design of Disco Moo and assigned a short descriptive code to each one. The codes concerned learning objectives, accessibility, feedback, engagement, and classroom participation. We then compared codes, discussed differences in interpretation, and grouped related codes into themes. We did not limit coding to our four interview questions. Five themes were especially relevant to the design.

\begin{enumerate}
    \item \textbf{Make computational thinking visible.} Panelists described block coding and Python as common approaches in middle- and high-school classrooms. One teacher emphasized helping students understand what happens ``under the hood'' rather than relying entirely on AI-generated code. Another noted that introductory games can teach the structure of instructions, including loops and conditionals. These comments shaped how we use tangible coding blocks. Reachy's movements are meant to make the consequences of a sequence observable, so the program is not left hidden inside generated code.

    \item \textbf{Scaffold the activity, not just the code.} Teachers cautioned against presenting middle-school students with a large, open-ended project all at once. One recommended breaking work into ``very distinct chunks''; another stressed clear verbal and visual prompts. We therefore organized the dance missions so that each one introduces a single challenge, with instructions available in more than one form.

    \item \textbf{Treat errors as opportunities to revise.} Panelists favored robot feedback that offers clues rather than the answer. As one put it, the robot should not respond with ``that's wrong, the answer is X.'' They also recommended keeping feedback specific and encouraging. We therefore moved away from ending a student's turn after an incorrect sequence. The proposed interaction uses robot prompts and interface cues that help students inspect, test, and revise their program.

    \item \textbf{Design participation into the rules.} Teachers warned that one student may dominate a shared game, particularly when they grasp the task more quickly than a partner. They suggested rotating roles and building turn-taking ``into the rules of your activity.'' We accordingly shifted from isolated individual turns toward collaborative planning, block assembly, testing, and debugging, with roles that can rotate across a mission.

    \item \textbf{Make classroom value and logistics explicit.} Panelists repeatedly asked whether a robotic game would strengthen the curriculum rather than merely attract attention. One teacher asked whether it would give the curriculum ``an extra bump'' or amount to ``pretty parts and pieces.'' They also raised preparation time, how many students would share a robot, and the need for ready-to-use materials. Disco Moo therefore has to make its programming objectives legible to teachers and include a practical plan for group rotation and for work away from the robot.
\end{enumerate}

Together, these themes suggest that an engaging robot is not, by itself, an educational activity. For Disco Moo, the design problem is to connect the appeal of robot movement to a visible programming sequence, a manageable progression of tasks, useful chances to revise, and a participation structure that works in a classroom. Our proposed interface will display the interpreted program and offer debugging cues, and Reachy Mini will give encouraging, hint-based responses. These are design intentions informed by the panel, not outcomes the panel established. Whether they improve learning or participation remains to be evaluated.

\subsection{Reflection}

The teacher panel helped us recognize that designing an educational robotic game involves more than selecting programming concepts. Although our original questions were useful for understanding classroom programming tools and learning difficulties, the discussions revealed additional considerations related to accessibility, student participation, and constructive feedback.

One important lesson was the distinction between making a robot entertaining and making it educational. While Disco Moo's cow-themed appearance and expressive dance movements may increase engagement, the gameplay must also provide opportunities to practice identifiable programming skills.

The panel also challenged our initial assumption that individual turn-taking would be an effective classroom format. Teachers highlighted the importance of keeping students engaged even when they are not directly controlling the robot, encouraging us to explore rotating roles and collaborative debugging.

In future interviews, we would ask more specific follow-up questions about appropriate challenge duration, group size, and methods for evaluating students' understanding of sequencing and loops. Presenting a storyboard or functional prototype could also help teachers provide more concrete feedback on the interaction design.



\section{Robot Concept and Design Process}

\subsection{Robot Role and Embodiment}

Disco Moo is a tangible programming game in which Reachy Mini acts as both a playful dance partner and an interactive coding companion. Inspired by Scratch, the game allows students to construct programs by physically assembling modular coding blocks instead of typing commands on a computer. Each block represents an instruction, such as turning the robot's head, nodding, or repeating a movement.

During gameplay, Reachy presents a dance challenge that students must solve by arranging the appropriate blocks in sequence. For example, Reachy might ask students to make it turn left, nod twice, and turn right. Students select and connect the corresponding blocks, present their program to the camera, and press Scan to start recognition and execution. Reachy's physical movements provide an observable representation of the programmed instructions, helping students understand sequencing and repetition.

Reachy is designed with a friendly, cow-inspired persona to make the interaction more playful and approachable for younger learners. We incorporated soft yarn around the robot's exterior to create a cow-like appearance, while retaining its expressive head movements. The robot's character is reflected in its dance-oriented challenges and planned encouraging feedback. Rather than functioning solely as a device that executes commands, Disco Moo is intended to act as a learning companion that encourages experimentation and collaborative problem-solving.

Based on the teacher panel feedback, we revised our initial individual turn-taking concept toward collaborative gameplay. Students can work together to plan sequences, assemble blocks, test their programs, and identify errors. The proposed interaction also includes progressively challenging missions and supportive feedback, allowing students to practice fundamental programming concepts without requiring prior Python experience.

\subsection{Iterative Design}

Our design process involved several stages of ideation, physical prototyping, and functional testing. We began by developing a Scratch-inspired game plan and identifying movement commands that could be represented through tangible coding blocks. Our initial concept used color-coded blocks to distinguish different command categories, including head movements, antenna movements, body rotation, and repetition.

As the design progressed, we developed a more consistent physical form inspired by Reachy Mini's appearance. We modeled rounded rectangular blocks with white bodies and black circular facial elements resembling Reachy's glasses. Four cylindrical studs on the top of each block allow multiple blocks to be stacked vertically, translating the concept of sequential programming into a physical assembly process.

We used 3D modeling and physical prototyping to investigate the blocks' connection mechanism. In particular, we produced variations with socket dimensions of 14.00 mm, 14.01 mm, 14.02 mm, and 14.03 mm to evaluate how dimensional tolerances affect the fit between blocks. This process allowed us to consider both connection stability and ease of assembly.

In parallel, we explored ArUco markers as a method for translating physical block arrangements into executable robot instructions. We printed the markers on paper and successfully tested their recognition and corresponding Reachy Mini movements through a camera-based interface. The interface includes a camera view and a Scan control for initiating execution. Although this established the feasibility of the recognition system, integrating the markers directly onto the 3D-printed blocks remains a future step.

The teacher panel also informed an important iteration in our interaction design. Our original game plan involved students completing challenges individually, with incorrect answers ending their turns. Following the teachers' recommendations, we shifted toward collaborative problem-solving, progressive challenges, and constructive feedback.

Throughout the process, we considered fabrication time, available materials, and technical feasibility. 3D printing allowed us to produce reusable, interlocking blocks with consistent geometry, while paper-based marker testing enabled us to evaluate the software before completing the physical integration. These decisions helped us develop a preliminary functional prototype within the lab timeframe while identifying areas for further refinement.



\section{Fabrication}

\subsection{Laser-Cut Nametag}

As part of the robot embodiment design, we created a custom laser-cut nametag for Disco Moo. The nametag serves as a visual identifier for Reachy's cow-inspired persona, reinforcing its identity as a playful coding companion.

We used laser cutting to produce the nametag because the process allows precise fabrication of flat shapes and engraved text. Compared with 3D printing, laser cutting is more suitable for this relatively simple, two-dimensional component and requires less fabrication time.

The completed nametag was incorporated into our robot embodiment alongside the soft yarn decorations used to create its cow-like appearance.

% TODO: Specify the nametag material and exact engraved text.
% TODO: Insert a photograph of the assembled nametag.

\subsection{Custom Laser-Cut Part}

Our fabrication process also explored how laser cutting could support the physical embodiment of Disco Moo. Laser cutting is particularly appropriate for flat decorative components that require precise outlines or engraved details, while 3D printing is better suited for the volumetric geometry of our interlocking coding blocks.

% TODO: Checkpoint 4 requires a custom laser-cut
% embodiment part, a photograph, the material used,
% a fabrication rationale, and the corresponding design file.
% Add the actual fabricated component here.

\subsection{3D-Printed Components}

We designed and fabricated modular coding blocks inspired by Reachy Mini's physical appearance. Each block features a rounded rectangular white body with two black circular elements connected by a narrow bridge, resembling Reachy's distinctive glasses. This visual design creates a consistent identity between the robot and the physical programming interface.

Each block incorporates four cylindrical studs on its upper surface, allowing multiple blocks to be stacked vertically. This interlocking mechanism enables students to physically construct a sequence of programming instructions, translating abstract coding structures into tangible arrangements.

During prototyping, we experimented with four socket dimensions: 14.00 mm, 14.01 mm, 14.02 mm, and 14.03 mm. These variations were designed to investigate how small dimensional differences influence the fit and stability of the connections. The tolerance testing formed an important part of our iterative fabrication process, as the blocks needed to remain connected while still being easy for students to assemble and separate.

We chose 3D printing because it allowed us to fabricate three-dimensional, reusable components with consistent shapes and integrated connection features. Unlike laser cutting or simple cardboard construction, 3D printing enabled the rounded geometry and cylindrical interlocking structures to be produced as part of the block design.

The physical blocks are intended to incorporate printed programming labels and ArUco markers. Each marker corresponds to a specific command that can be recognized through the camera-based interface and translated into a robot movement. We successfully tested marker recognition and robot execution using paper-printed ArUco markers. However, the markers have not yet been integrated and tested directly on the 3D-printed blocks.

The resulting prototypes demonstrate the physical feasibility of the modular coding system and provide a foundation for integrating tangible programming with Reachy Mini's movements.

% TODO: Insert photograph of the assembled 3D-printed blocks.
% TODO: Include corresponding 3D modeling and printing files.
% TODO: Specify the selected socket dimension and test outcome.




\section{Final Robot Design}

Our final Lab 4 prototype of Disco Moo combines a cow-inspired Reachy Mini embodiment with a tangible block-based programming system. The robot is decorated with soft yarn to establish a playful, approachable character, while the modular coding blocks adopt Reachy's recognizable black-and-white appearance. Together, these elements create a consistent visual identity between the robot and the physical programming interface.

The design evolved through sketching and conceptual planning, 3D modeling, physical prototyping, and functional testing. Our initial concept focused on individual coding challenges using color-coded blocks. Following the teacher panel, we revised the gameplay toward collaborative problem-solving, progressive learning, and constructive feedback. We also developed a consistent physical design for the blocks, incorporating Reachy-inspired facial details and interlocking structures.

The fabrication process required balancing usability, material constraints, and production time. We used 3D printing to produce reusable coding blocks with rounded geometry and integrated connection features. Testing different socket dimensions allowed us to investigate the stability and usability of the interlocking mechanism. Meanwhile, laser cutting provided an efficient method for producing the robot's personalized nametag, and yarn offered a lightweight, low-cost approach to modifying Reachy's appearance without requiring a fully redesigned body.

From an affordance perspective, the cylindrical studs and corresponding sockets communicate how the blocks can be physically connected. Their stackable structure encourages students to arrange instructions in a linear sequence, making the relationship between program order and robot behavior more tangible. The rounded shapes and recognizable facial features also help distinguish the blocks as interactive learning objects rather than conventional computer components.

Materiality plays an important role in the interaction experience. The rigid 3D-printed blocks provide a physical representation of programming commands, while the soft yarn surrounding Reachy creates a contrasting, approachable appearance. This combination of materials supports the game's playful identity and reinforces the relationship between physical manipulation and digital execution.

The interaction design connects three elements: the student, the physical coding blocks, and Reachy Mini. Students construct programs by selecting and stacking blocks, present the sequence to the camera, and press Scan to start execution. Reachy's resulting movements provide immediate, observable feedback on the interpreted commands. This interaction transforms programming from an exclusively screen-based activity into a tangible and embodied learning experience.

For our selected fabricated component, we highlight the 3D-printed modular coding block. Its interlocking geometry enables physical program construction, while its Reachy-inspired appearance visually connects the programming interface to the robot. Although the blocks function as external interaction objects rather than attachments to Reachy's body, they contribute to the overall embodiment and interaction concept of Disco Moo.

Our functional testing demonstrated that paper-printed ArUco markers could be recognized and translated into Reachy Mini movements. However, the complete integration of markers onto the 3D-printed blocks has not yet been tested. Future iterations will focus on integrating the physical and digital components, refining block recognition, and evaluating the interaction with students.

\begin{figure}[t]
    \centering
    \figureplaceholder{Insert photo of the final Disco Moo prototype.}
    \caption{Final Lab 4 prototype of Disco Moo, combining the cow-inspired Reachy Mini embodiment with modular 3D-printed coding blocks.}
    \label{fig:finalrobot}
\end{figure}


\section{Discussion and Future Work}

\subsection{Additional Research Output}

In addition to the physical coding blocks and the robot embodiment, we developed a camera-based interface that reads an ordered stack of ArUco markers and sends the corresponding commands to Reachy Mini. This is the screen a student would use to check which program the camera recognized and which step Reachy is carrying out. The prototype is titled Reachy LEGO. Figures~\ref{fig:ui-states} and~\ref{fig:ui-execution} show it during a feasibility test with paper markers.

The screen has four parts: a status bar, a short instruction, a live camera view, and a Sequence panel. In the ready state, the status bar reads READY and reports that the interface is connected to Reachy. The instruction asks the student to hold the stack in front of Reachy and press Scan. Until that scan succeeds, the Sequence panel states that the scanned sequence will appear there, as in Fig.~\ref{fig:ui-ready}.

Pressing Scan moves the interface to RUNNING. Each detected marker is outlined and labeled in the camera view. The Sequence panel then lists the moves in order, placing a plain-language name beside the command sent to the robot, and highlights the step currently underway. In Fig.~\ref{fig:ui-running}, the recognized program is twist body right, look right, and shake head. The status line names the first of those moves while Reachy performs it, and the Scan control changes to Please wait so the stack is not read again mid-dance.

\begin{figure*}[t]
    \centering
    \subfloat[Ready state. The student holds a vertical stack of ArUco markers in view of the camera. The interface asks them to press Scan, and the sequence list is still empty.]{%
        \includegraphics[width=0.48\textwidth]{"additional research output/custom UI - hold up markers.jpg"}%
        \label{fig:ui-ready}}
    \hfil
    \subfloat[Running state. Detected markers are labeled on the camera view, and the sequence panel highlights twist body right as the move Reachy is performing.]{%
        \includegraphics[width=0.48\textwidth]{"additional research output/custom UI - scanning in process.jpg"}%
        \label{fig:ui-running}}
    \caption{Custom interface a student would see while presenting a paper marker stack to Reachy Mini.}
    \label{fig:ui-states}
\end{figure*}

The same display stays beside the robot while the sequence runs, so the highlighted line and the physical movement can be checked against each other. Fig.~\ref{fig:reachy-exec} shows Reachy Mini during that test. The interface is RUNNING a six-step sequence and highlights twist body right as the active step. The camera view shows the marker sheet used for that longer program, with each detected marker labeled.

When the last move finishes, the status changes to CELEBRATING and the interface reports that the dance is complete. The full program remains on screen, including the action keys for the recognized markers, as in Fig.~\ref{fig:ui-complete}.

\begin{figure*}[t]
    \centering
    \subfloat[Reachy Mini executing the highlighted step. The interface lists a six-step sequence and marks twist body right as the move in progress.]{%
        \includegraphics[width=0.48\textwidth]{"additional research output/Reachy execute order and indicate which steps .jpg"}%
        \label{fig:reachy-exec}}
    \hfil
    \subfloat[Complete state. The status bar reads CELEBRATING, the interface reports that the dance is complete, and the full three-step sequence remains visible.]{%
        \includegraphics[width=0.48\textwidth]{"additional research output/complete state.jpg"}%
        \label{fig:ui-complete}}
    \caption{Reachy Mini carrying out a scanned sequence while the interface indicates the active step, and the screen after the dance finishes.}
    \label{fig:ui-execution}
\end{figure*}

These trials show that the interface can read an ordered stack, display the interpreted program, and have Reachy execute the listed moves one step at a time. The markers in this test were printed on paper. Attaching them to the 3D-printed blocks, and keeping each marker visible to the camera once those blocks are stacked, remains future work. For the selected fabricated component, we retain the 3D-printed modular coding block described in Section 5. Its interlocking structure and recognizable visual design provide the physical foundation for the proposed programming interaction.



\subsection{Reflection}

Developing Disco Moo helped us understand how physical fabrication, interaction design, and educational objectives must work together in human-robot interaction. One of our main challenges was translating abstract programming concepts into physical objects that are both intuitive to assemble and technically recognizable.

The iterative fabrication process highlighted the importance of small dimensional adjustments in 3D printing. Testing different socket tolerances demonstrated how physical connection design can influence the usability of modular learning tools. Meanwhile, the paper-based ArUco tests allowed us to verify the basic software-to-robot interaction before completing the physical integration.

The teacher panel also significantly influenced our design approach. Initially, we focused on creating an entertaining individual coding challenge. However, the teachers' feedback encouraged us to consider scaffolding, accessibility, collaboration, and constructive feedback as essential components of an educational experience. This shifted our focus from simply making Reachy perform movements toward helping students understand the programming logic behind those movements.

Despite these developments, the current prototype has several limitations. Although ArUco recognition and robot execution were successfully tested using paper markers, the complete system has not yet been evaluated with markers attached to the 3D-printed blocks. Additionally, the proposed collaborative gameplay, progressive challenges, and debugging feedback require further implementation and testing.

For the final project, we plan to integrate ArUco markers and programming labels onto the physical blocks, refine the recognition of assembled command sequences, and further develop the interactive game interface. We also aim to evaluate whether students can understand sequencing, repetition, and debugging through the tangible coding experience. These next steps will help us move from a preliminary functional prototype toward a more complete educational robotic game.


\end{document}
